\subsection{Sprint 0}

A Sprint 0 é projetada para ser uma espécie de aquecimento para as próximas sprints. Durante essa fase, não se desenvolve ou se desenvolve muito pouco da aplicação; em vez disso, o foco está na definição das tecnologias, user stories, workflows e outros aspectos essenciais. No entanto, apesar da aparência inicial de tranquilidade, esta Sprint é longe de ser simples, exigindo estudo e dedicação intensos para garantir que todos os membros da equipe estejam bem preparados para o projeto e possam progredir mais suavemente nas próximas sprints.

Ao longo das reuniões desta sprint, nossa equipe concentrou-se na definição das estruturas básicas do projeto. Isso incluiu a escolha das tecnologias a serem utilizadas, a configuração do banco de dados, a identificação das User Stories, criação de mockups iniciais e outros aspectos fundamentais.

Ao contrário da AGES anterior, minha experiencia adquirida anteriormente me fez não passar por tantas dificuldades estudando novos tópicos da aplicação. A verdadeira novidade do projeto para mim seria em relação a blockchain em sí. Durante a primeira sprint meu foco seria entender de fato o que seria ela e como funcionava, suas vantagens e desvantagens e como implementaríamos em nossa aplicação back-end. Além da blockchain, estudei o básico de spring boot apenas para me familiarizar com a técnologia. Em relação aos mockups não me envolvi em suas criações por não ter habilidade alguma com a criação deles. Por fim comecei o desenvolvimento da modelagem conceitual do nosso banco de dados, no entanto, com a escassez de users stories não pude ir muito a fundo com a modelagem.

A segunda reunião da sprint com o stakeholder foi significativamente mais curta do que a primeira. Durante esse encontro, apresentamos os protótipos de telas (mockups) e todas as user stories desenvolvidas. Os AGES IV decidiram que a próxima sprint seria a mais desafiadora até então, selecionando 5 das 12 user stories para serem concluídas. No final da reunião, ficamos felizes ao ver que o stakeholder estava muito satisfeito com o progresso apresentado pelo time e não solicitou nenhuma alteração, o que foi uma grande vitória para todos nós.

Apesar do grande êxito demonstrado durante a reunião com o stakeholder, esta sprint também teve seus momentos desafiadores que impactariam toda a sequência do semestre. Durante a sprint, dois dos quatro AGES II saíram do projeto, deixando a responsabilidade do banco de dados para mim e outro colega. Apesar da situação, tive confiança de que nós dois seríamos capazes de lidar com o trabalho. No entanto, essa mudança abalou um pouco a equipe de forma geral.